\section{Limitations}
\label{sec:limitations}

This study has several important limitations that constrain the generalizability and interpretability of its findings.

\subsubsection*{Model scope.} Primary experiments were conducted with Llama-3.1-8B-Instant. While cross-model validation was performed with Llama-3.3-70B, this was limited to two conditions (C0, C3) and a single inference run, precluding a full nine-condition multi-seed ablation for the larger architecture. The observed effects—including the task-dependent performance asymmetry—may be sensitive to model family, parameter count, or inference backend beyond the two evaluated scales.

\subsubsection*{Single-split evaluation.} Results are based on a single scaffold-split test set for each dataset (BBBP $N=204$, BACE $N=152$) without full cross-validation. This makes the reported ROC-AUC values sensitive to the specific test partition. While we employed multi-seed resampling to evaluate generative stability (Table~\ref{tab:robustness}), extending this framework to multiple independent scaffold splits is required to definitively isolate generalizable phenomena from dataset-specific artifacts. We explicitly defer this multi-split validation to future work, positioning the current findings as exploratory behavioral probes rather than comprehensive benchmarks.

\subsubsection*{Limited sample size.} With 204 and 152 molecules respectively, this study is underpowered for detecting small but potentially meaningful effect size differences between conditions. The SP category comprised only 5 naturally occurring instances in BBBP, making any distributional claims about SP highly tentative.

\subsubsection*{C2 control confound.} The chemical priming condition (C2) contains domain-relevant vocabulary (``structural'', ``molecular'', ``physicochemical''), making it a partial rather than pure stylistic control. While C2b addresses this with non-chemical vocabulary, the C2 results should be interpreted with this confound in mind.

\subsubsection*{Taxonomy and Classifier Validity.} The heuristic classification scheme used to categorize free hallucinations (C3) into SP, PI, MF, or CC is an unvalidated rule-based system. To assess its reliability, we conducted an AI-assisted semantic validation on a stratified sample ($N=70$) using an LLM-as-judge approach (Llama-3.1-8B-Instant). The agreement between the heuristic classifier and the semantic judge was poor (Cohen's $\kappa = 0.117$, raw agreement 44.3\%), with the LLM judge exhibiting severe overclassification toward Mechanism Fabrication (MF) and failing to reliably identify Structural Phantoms (SP) or Property Inversions (PI). This high disagreement rate highlights the inherent ambiguity in categorizing generative chemical text. Consequently, the categories (SP, PI, MF, CC) should be interpreted strictly as \textit{prompt design intents} or heuristic keyword clusters, rather than definitive semantic truths. The Contextual Confabulation (CC) category functions as a catch-all residual class, encompassing heterogeneous subtypes whose homogeneous treatment obscures important variations. This limitation is further evidenced by the post-hoc verification of C4a--c outputs (Table~\ref{tab:c4_compliance}): despite prompt-constrained forcing, C4a achieved 0\% target compliance (SP), C4b achieved 25.5\% (PI), and C4c achieved 20.4\% (MF), with CC dominating all three conditions ($\geq$74\%). This systematic prompt--output misalignment demonstrates the challenge of enforcing strict semantic perturbation categories in generative chemistry pipelines.


\subsubsection*{Data leakage risk.} As highlighted in recent literature~\cite{guo2023benchmark, white2023assessment}, modern LLM pre-training corpora (such as those used for Llama and Gemma) likely encompass massive amounts of public chemical data, including the MoleculeNet datasets. Consequently, the zero-shot predictions evaluated here may partially reflect varying degrees of memorized retrieval rather than pure inferential reasoning. While our control conditions (C5) demonstrate that structural prompts without semantic alignment fail to produce predictive gains, hallucinated text might inadvertently trigger memorized associations, collapsing the evaluation framework from a test of reasoning to a test of prompted recall.

\subsubsection*{API non-determinism.} Despite setting temperature $= 0.0$ and seed $= 42$, Groq's API documentation notes that seed-based reproducibility is provided on a best-effort basis. Exact numerical reproducibility across runs is therefore not guaranteed.

\subsubsection*{ESOL Task Framing.} Our ESOL evaluation used a zero-shot regression prompt. However, as LLMs are typically optimized for categorical text output rather than high-precision scalar regression, the reported RMSE reflects the inherent limitations of zero-shot regression in the absence of specialized prompting or fine-tuning. These values should be interpreted as a relative measure of semantic sensitivity across augmentation conditions rather than an absolute measure of solubility prediction accuracy.

\subsubsection*{Dataset scope.}
The primary nine-condition ablation was conducted on two 
MoleculeNet benchmarks (BBBP and BACE), with targeted 
three-condition generalization validation on Tox21 and 
ESOL. Full nine-condition ablation on Tox21 and ESOL, 
and extension to additional task types such as protein 
binding affinity, remains reserved for future work. 
Cross-dataset generalization of the observed 
distributional compression phenomenon should be 
interpreted with caution given the limited condition 
coverage of the extended benchmarks.

\subsubsection*{Inference-only evaluation.} While we include a traditional ML baseline (Random Forest) for context, we evaluate only zero-shot LLM inference. Comparison with fine-tuned LLM models would provide additional context for the practical significance of the observed effects.

\subsubsection*{Methodological confounds.} The BACE Mechanism Fabrication (C4c) condition directly prompts the LLM to hypothesize biological targets. Because BACE-1 is the implied target of the BACE dataset, the MF text may inadvertently contain information about BACE-1 itself, creating a direct semantic confound rather than a general structural perturbation. Although we validated the C5 control across multiple permutation seeds ($0.49 \pm 0.02$), a larger number of permutations with full-dataset evaluation would further strengthen this control.
